A single intraclass correlation (ICC) hides how event construction shapes a binocular saccade marker's repeatability. In an exploratory, post hoc analysis of the GazeBaseVR reading task (407 participants, 250\,Hz head-mounted display), we contrasted left/right-eye peak-velocity ratios between return sweeps and progressive saccades under six event constructions. In 376 participants eligible under all six, admitting events detected on the mean of both eyes' speeds only when a matched per-eye detection overlapped them raised the absolute-agreement ICC(A,1) between two sessions about 30 minutes apart from .374 to .693 (difference .319, descriptive 95\% CI .003 to .501). One participant's session with implausibly fast one-eye events carried most of this; omitting them, it was .054 ($-$.007 to .129). A conventional 1,000$^\circ$/s peak-velocity ceiling gave .711, and across detector settings the unfiltered ICC ranged from .226 to .703 while admitted and ceiling ICCs stayed at .66--.74. Recording eligibility, by excluding two participants, also moved the ICC substantially. Among simulated one-eye faults reaching it, recording-centred abstention rejected nearly all 4$^\circ$ and 8$^\circ$ steps but fell short of its frozen criterion for threefold slowings (77\%) and 6$^\circ$ peak glitches (38\%). Reliability reports for binocular kinematic markers should state event-plausibility and eligibility rules and single-participant influence.
